\cventry{Feb 2021 - \textbf{Present}}{{R\&D}}{\textsc{\href{https://www.stiima.cnr.it/}{\color{blue}STIIMA}}}{Lecco}{}
{Field: \textbf{MedTech}
\vspace{0.1cm}
%\begin{itemize}
%\item Identifying and interpreting patient and clinical need with a focus on uncovering insights for design innovation;
%\item Developing and building physical prototypes based on what users need, aligning with business goals and improving user experiences;
%\item Expertise in parametric 3D CAD software such as Solidworks and Fusion 360;
%\item Familiarity with ISO as ISO 13485, ISO 9001, 10993 for CE mark, technical files and dossiers;
%\item Integration prototypes into demonstration and testing setups;
%\item Propose coherent manufacturing technique, designing for manufacturing purpose according to the batch size;
%\item Strong collaboration skills with cross-functional teams and the ability to operate autonomously in a dynamic environment;
%\end{itemize}}
%\begin{itemize}
%\item Strong problem-solving skills and ability to collaborate with cross-functional teams, working autonomously in dynamic environments;
%\item Demonstrating creativity and innovation across both physical and digital design, developing solutions that meet patient and clinical needs;
%\item Creating and iterating UX and design artifacts, including research summaries, sketches, user scenarios, product renderings, and packaging designs;
%\item Proficiency in parametric 3D CAD software, including SolidWorks and Fusion 360, with knowledge of PDM systems;
%\item Designed, fabricated and assembled over 10 functional prototypes;
%\item Integrating 3D printed prototypes into demonstration and testing setups, supported by performance and usability testing;
%\item Designing for manufacturability by proposing suitable techniques based on batch size, ensuring efficiency and cost-effectiveness;
%\item Familiarity with standards such as IEC 60601, ISO 13485, ISO 9001, and ISO 10993 for CE marking, including creating technical files, operating manual and dossiers;
% \item Experience with SAAS product design and understanding of the designer's role within the supply chain;
%\item Collaborated with clinicians, engineers, and researchers to design and test MedTech rehabilitation prototypes.
%\item Designed, fabricated, and assembled 10+ functional prototypes using SolidWorks, Fusion 360, and 3D printing.
%\item Produced UX and design artifacts including research summaries, sketches, user scenarios, product renderings, and packaging concepts.
%\item Supported CE-oriented documentation, including technical files, operating manuals, and design dossiers.
%\item Proposed manufacturing methods based on batch size, cost, usability, and technical constraints.

%LAST USED
%\item Collaborated with clinicians, engineers, and researchers to design, prototype, and test sensor-based MedTech rehabilitation systems.
%\item Designed, fabricated, and assembled 10+ functional prototypes using SolidWorks, Fusion 360, and 3D printing.
%\item Supported device development through usability testing, technical validation, and iterative design improvements.
%\item Contributed to CE-oriented documentation, including technical files, operating manuals, design dossiers, and verification records.
%\item Applied risk-based thinking to design decisions, considering usability, safety, manufacturability, cost, and technical constraints.
    %\item \textbf{Project V.I.T.A.}:It is a digital platform using VR and AR to support individuals with autism through immersive therapy, play, and socialization, with remote support and activity tracking.
                %\begin{itemize}
                %\item Role: Development of an IoT distributed and coordinated system which provide assistive feedback, such as light, sound, or vibration patterns delivered at different frequency and intensity, useful to the patient for space orientation training;
                %\item Tasks: Defining the problem, setting objectives, and outlining key requirements. It includes system design and prototyping through component selection, assembly, software integration, BLE pairing, sensor(s) calibration, sync, battery, device management, and stability. Performance is evaluated through testing, which involves designing test cases, collecting data, and analyzing results;
                %\end{itemize}  
\vspace{-1em}
\cventry
{\raisebox{-0.60\height}{\includegraphics[height=1cm]{Pictures/cnr.png}}}
{{R\&D Product Engineer}}{\textsc{\href{https://www.stiima.cnr.it/}{\color{blue}STIIMA}}}{Lecco}{}
{%Field: \textbf{MedTech}}
Feb 2021 - Present}
\vspace{0.1cm}
\begin{itemize}
\item Developed sensor-based rehabilitation and assistive technology systems, translating clinical requirements and user needs into functional prototypes for experimental validation.
\vspace{+0.1em}
\item Defined the architecture of multimodal IOT feedback platforms, integrating wearable sensing, indoor localization, BLE communication, and adaptive audio, visual, and haptic feedback.
\vspace{+0.1em}
\item Delivered 4+ functional prototypes across embedded electronics, wearable devices, and interactive systems, supporting projects from concept development to technology demonstration.
\vspace{+0.1em}
\item Designed communication infrastructures and evaluation protocols for digital health applications, assessing usability, system performance, and user interaction through experimental testing.
\end{itemize}
\vspace{+0.1em}
\texttt{Embedded Systems • IoT • ESP32 • Arduino • PlatformIO • VS • BLE GATT • ROS2 • C++ • C\# • Python • Prototyping • CAD • FDM • Material Selection}

